\section{有限位数楕円境界を持つ四係数族の中心上有限な相対模型}
\label{sec:four-parameter-central-model}

この節では、固定した周囲の代数 $B$ における
中心三次斉次元 $G$ の商だけを動かす。
係数環と族を
\[
 P=\mathbb C[a,b,c,d],\qquad
 R=P[u,v],\qquad
 \mathcal A=B\otimes_{\mathbb C}P/
 (G+au^3+bu^2v+cuv^2+dv^3)
\]
と書く。$a,b,c,d$ の通常の次数は零、$u,v,x,y,z$ の通常の次数は一とする。
以下で用いる入力は、既存の PBW 構成により
\begin{equation}\label{eq:central-model-filtered-input}
 \operatorname{gr}_{x,y,z}\mathcal A
 \simeq H\otimes_{\mathbb C}P[u,v],\qquad
 H=B(E,\mathcal L,\tau)
\end{equation}
となり、$E$ が滑らかな楕円曲線、$\mathcal L$ が豊富な直線束、
$\tau$ が有限位数の平行移動であることである。
各 $F_j\mathcal A$ は $R$ 上有限自由で、閉点への基底変更と filtration が両立する。
また $\mathcal A$ は通常の次数一で生成され、$\mathcal A_0=P$ とする。
具体的な $B_*$ については、命題\ref{prop:moduli-rigid-example}の
PBW条件と楕円境界の同定を入力とする。
境界の平行移動の位数六は補題\ref{CSBGlem:order-six}で確認される。

式 \eqref{eq:central-model-filtered-input} の各斉次部分は $R$ 上自由であるから、
その基底を filtration に沿って持ち上げると
\begin{equation}\label{eq:central-model-R-free}
 \mathcal A\text{ は次数付き }P[u,v]\text{-加群として自由}
\end{equation}
である。
特に任意の可換 $P$-代数 $S$ について、
$\mathcal A_S=\mathcal A\otimes_P S$ は $S[u,v]$ 上自由であり、
$u,v$ およびその正冪は普遍的な中心正則列となる。

\subsection{中心元の一様な持上げ}

補題\ref{CSBGlem:global-center-lifting}の中心持上げを，
$x,y,z$に関するfiltrationへ適用する。

\begin{theorem}[全四係数族の有限な中心模型]
\label{thm:central-model-finite}
上述の入力の下で、有限型次数付き可換 $P$-代数
\[
 P[u,v]\subset C\subset Z(\mathcal A),\qquad C_0=P
\]
であって $\mathcal A$ が $C$ 上有限となるものが存在する。
任意の noetherian 可換 $P$-代数 $S$ に対し、
$\mathcal A_S$ は $C_S=C\otimes_P S$ 上有限であり、左右 noetherian である。
\end{theorem}

\begin{proof}
補題\ref{CSBGlem:global-center-lifting}を $R_0=P[u,v]$、$U=\mathcal A$ に適用する。
そこで持ち上げる中心元 $z\in Z(H)_j$ は、通常の次数も $j$ である。
中心に持ち上げた元を通常の次数で分解すると、各斉次成分も中心に入る。
通常の次数 $j$ の成分を取れば、最高 filtration symbol は依然として $z$ である。
従って中心生成元の持上げを通常の次数で斉次に選べる。
$P[u,v]$ と有限個のこれらの持上げが生成する代数を $C$ とする。
同補題の filtration 帰納法は $\mathcal A$ の $C$ 上の有限性を与える。

有限生成加群であることは任意の基底変更後にも保たれる。
$C_S$ は noetherian 環 $S$ 上有限型の可換代数なので noetherian である。
$\mathcal A_S$ の左イデアルまたは右イデアルの昇鎖は $C_S$-部分加群の昇鎖であり、停止する。
\end{proof}

\begin{remark}[全中心および非被約基底の扱い]
\label{rem:central-model-not-full-center}
ここでは $C=Z(\mathcal A)$ を主張していない。
また、任意の $P\to S$ の後に $C_S\to\mathcal A_S$ が単射であるとも主張しない。
必要なのは、中心に像を持つ環準同型
$C_S\to Z(\mathcal A_S)$ と、その作用に関する有限加群性だけである。
中心持上げに Rains の定理と被約性を使うのは $P[u,v]$ 上で一度だけであり、
非被約な $S$ 上でそれらを再適用していない。
\end{remark}

\begin{corollary}[剛体族の普遍的な二次中心]
\label{cor:central-model-universal-center-two}
命題\ref{prop:moduli-rigid-example}の$B_*$とその四係数族について，
任意の可換$P$-代数$S$に対し
\[
 Z(\mathcal A_S)_2=S u^2\oplus S uv\oplus S v^2
\]
である．従ってこの族は定義\ref{def:abel-two-center-family}の
普遍的な中心二変数条件を全係数基底上で満たす．
この主張では$S$のnoetherian性を仮定しない．
\end{corollary}
\begin{proof}
三次商は二次関係を変えないので，$(\mathcal A)_2$は
$(B_*)_2\otimes P$であり，$Pu^2\oplus Puv\oplus Pv^2$は
その定数な直和因子である．従って商は自由rank~12である．
$(\mathcal A)_1$は自由rank~5，$(\mathcal A)_3$は自由rank~34だから，
命題\ref{prop:abel-center-rank-open}の有限自由加群間の写像
\[
 (\mathcal A)_2/\operatorname{Sym}^2(Pu\oplus Pv)
 \longrightarrow\operatorname{Hom}_P((\mathcal A)_1,(\mathcal A)_3)
\]
を使える．各閉点$s$では，三次元$G_*+f_s$の最高symbolは非零正則な$\gamma$であり，
補題\ref{lem:moduli-low-center}と\ref{lem:moduli-no-new-quadratic-center}から
$Z(\mathcal A_s)_2=k(s)u^2\oplus k(s)uv\oplus k(s)v^2$である．
ここで固定した主変数の三次式の滑らかさと境界translationの位数六を使う．
従って全閉点がrank~12の開部分に入る．$P$はJacobson環なのでその開部分は全体である．
写像は局所的に分裂単射であり，任意の$S$への係数変更後も単射である．
全次数の有限自由性と$S[u,v]$-自由性も既に示したので，定義の他の条件も満たす．
\end{proof}

\subsection{射影スキーム上の有限代数層と普遍的な section 復元}

以下、$\ell>0$ を十分割り切れる偶数とし、$D=C^{(\ell)}$ が $D_1$ で生成されるように取る。
\[
 Y=\operatorname{Proj}D,\qquad
 \mathscr A=\widetilde{\mathcal A^{(\ell)}}
\]
と置く。$Y\to\operatorname{Spec}P$ は射影的で、$\mathscr A$ は有限な
$\mathcal O_Y$-代数である。任意の可換 $P$-代数 $S$ について
$Y_S=\operatorname{Proj}(D\otimes_P S)$ と書き、$\mathscr A_S$ を基底変更した代数層とする。
中心準同型に核が生じる場合も、この模型をその support とともに用いる。

\begin{proposition}[相対 tails の有限代数層模型]
\label{prop:central-model-qgr}
任意の可換 $P$-代数 $S$ について、基底変更に自然な同値
\[
 \operatorname{QGr}(\mathcal A_S)
 \simeq \operatorname{QCoh}(Y_S,\mathscr A_S)
\]
がある。$S$ が noetherian なら、これは同値
\[
 \operatorname{qgr}(\mathcal A_S)
 \simeq \operatorname{coh}(Y_S,\mathscr A_S)
\]
に制限される。
構造対象 $\mathcal O=\pi\mathcal A_S$ は $\mathscr A_S$ に対応し、
$\mathcal O(m)$ は
$\widetilde{\mathcal A_S(m)^{(\ell)}}$ に対応する。
大きい加群圏は $S$-線型 Grothendieck 圏であり、
$S$ が noetherian なら locally noetherian である。
これらの同値は任意の基底変更の合成および fppf 降下に両立する。
\end{proposition}

\begin{proof}
標準開 $D_+(c)$、$c\in D_1=C_\ell$ を取る。
$\mathcal A_S$ は次数一生成なので、$c$ を次数一元の長さ $\ell$ の積の和で書ける。
中心局所化 $T=\mathcal A_S[c^{-1}]$ において、その式を最初または最後の因子で切れば
\[
 1\in T_1T_{-1},\qquad 1\in T_{-1}T_1
\]
となる。
従って $T$ は強く $\mathbb Z$-次数付きであり、次数零部分を取る関手は
次数付き $T$-加群と $T_0$-加群の同値を与える。
実際 $T_n$ は可逆な $T_0$-両側加群であり、
$N\mapsto N\otimes_{T_0}T$ が逆関手である。

一方、$\mathcal A_S$ が $C_S$ 上有限なので、$\mathcal A_S$-torsion と
$C_{S,+}$-torsion は一致する。
例えば $mC_{S,\geq n}=0$ なら、有限個の斉次 $C_S$-加群生成元の次数の最大値を $e$ として
$m\mathcal A_{S,\geq n+e}=0$ となる。
逆方向は $C_S$ の像が中心にあることから直ちに従う。
標準開上の上記の同値を貼り合わせると tails の同値を得る。
これは通常の Serre の Proj 対応の証明を有限な代数層に適用したものであり、
中心局所化と次数零部分だけで構成される。
逆向きには、quasi-coherent な $\mathscr A_S$-加群 $N$ から
\[
 \bigoplus_{n\in\mathbb Z}
 \Gamma(Y_S,N\otimes_{\mathcal O_{Y_S}}\mathcal O_{Y_S}(n))
\]
という次数付き $\mathcal A_S^{(\ell)}$-加群を取る。
この加群の層化は $N$ を復元する
\cite[Lemma~28.29.5, Tag~0AG5]{StacksProj}。
さらに次数を $\ell$ 倍にして $\mathcal A_S$ へ誘導すれば、
上記の強次数付き局所化の逆関手を各標準開上で得る。
有限中心被覆上で層化の核が torsion と一致することから、
これらは商圏上で互いに逆である。
ここまでの構成は $S$ の noetherian 性を使わない。
また中心局所化と次数零部分は任意の係数拡大と可換する。
したがって基底変更とその合成に自然である。
coherent 部分の対応には $S$ の noetherian 性を用いる。

有限代数層の quasi-coherent 加群は通常の加群とその作用写像の降下として
fppf 降下を持つ。大きい tails 商は Grothendieck 圏であり、
$S$ が noetherian なら $\mathcal A_S$ も noetherian なので
この商圏は locally noetherian である。
固定した二つの模型の間の同値および自然変換の降下には，
補題\ref{lem:abel-canonical-family-equivalences}を用いる。
\end{proof}

\begin{theorem}[任意基底変更後の全次数 saturation]
\label{thm:central-model-universal-sections}
任意の可換 $P$-代数 $S$ と任意の整数 $m$ に対し、標準写像
\[
 (\mathcal A_S)_m\longrightarrow
 \operatorname{Hom}_{\operatorname{QGr}(\mathcal A_S)}
 (\mathcal O,\mathcal O(m))
\]
は同型である。これらの同型は任意の可換な基底変更、乗法および合成と両立する。
特に負次数の右辺も零である。
非noetherianな$S$については大きい$\operatorname{QGr}$を用いる。
\end{theorem}

\begin{proof}
まず$S$がnoetherianである場合を示す。
$D_S=D\otimes_P S$、$I=D_{S,+}$ とし、次数付き有限 $D_S$-加群
$M_m=\mathcal A_S(m)^{(\ell)}$ を考える。
式 \eqref{eq:central-model-R-free} により、$u^\ell,v^\ell\in (D_S)_1$
は $M_m$ 上の正則列である。
商を取ることと Veronese 部分を取ることは可換なので、
この主張は $\mathcal A_S$ が $S[u,v]$ 上自由であることから従う。

$u^\ell$ の単射性から $H^0_I(M_m)=0$ である。
また $M_m/u^\ell M_m$ 上では $v^\ell$ が単射なので
$H^0_I(M_m/u^\ell M_m)=0$ となる。
local cohomology の長完全列により、$u^\ell$ は $H^1_I(M_m)$ 上でも単射である。
しかし $H^1_I(M_m)$ は $I$-torsion であり、各元は $u^\ell$ のある冪で消える。
従って $H^1_I(M_m)=0$ である。

斉次生成元による増大\v{C}ech複体とその次数成分を比較して得られる，
標準次数付き可換環のProjに対するlocal cohomologyの完全列
\[
 0\longrightarrow H^0_I(M_m)\longrightarrow M_m
 \longrightarrow\bigoplus_{n\in\mathbb Z}
 H^0(Y_S,\widetilde M_m(n))\longrightarrow H^1_I(M_m)
 \longrightarrow0
\]
の次数零部分を取る。
$\widetilde M_m=\mathcal O(m)$、$\mathcal O=\mathscr A_S$ なので、
\[
 H^0(Y_S,\widetilde M_m)
 =\operatorname{Hom}_{\mathscr A_S}(\mathscr A_S,\widetilde M_m)
\]
であり、所望の同型を得る。
ここで用いた写像は、元を左乗法による加群準同型へ送る標準写像そのものである。
この段階では、同じ証明が全てのnoetherianな$S$に適用できる。

一般の$S$を、その有限型$P$-部分代数$S_\lambda$のfiltered colimitとして書く。
$\mathcal D=\operatorname{QGr}(\mathcal A)$とし、$E_m=\mathcal O(m)$と置く。
$\mathcal D$はlocally noetherianであり、各$E_m$は有限表示対象である。
補題\ref{lem:abel-arbitrary-coefficient-tails}により
$\operatorname{QGr}(\mathcal A_S)\simeq\mathcal D_S$であり、
自由係数作用と忘却の随伴から
\begin{align*}
 \operatorname{Hom}_{\mathcal D_S}(S\otimes_P E_0,S\otimes_P E_m)
 &\simeq\operatorname{Hom}_{\mathcal D}(E_0,S\otimes_P E_m)\\
 &\simeq\varinjlim_\lambda
       \operatorname{Hom}_{\mathcal D}(E_0,S_\lambda\otimes_P E_m)\\
 &\simeq\varinjlim_\lambda(S_\lambda\otimes_P\mathcal A_m)
  =S\otimes_P\mathcal A_m
\end{align*}
となる。各$S_\lambda$はnoetherianなので、最後から二番目の同型には既に証明した場合を使える。
これは標準section写像を計算したものであり、基底変更前後の正方形は可換である。
従って同型は任意の基底変更、乗法および合成と両立する。
\end{proof}

\subsection{偶数列に関する補足}

次の各主張は、上の saturation と区別して記載する。
相対Serre双対は第\ref{sec:relative-serre-projective}節で別に構成する。

\begin{proposition}[偶数列の豊富性、復元、関手の決定]
\label{prop:central-model-even-sequence}
noetherian な可換 $P$-代数 $S$ を固定し、
上記の模型において $E_i=\mathcal O(-2i)$、$i\geq0$ とする。
\begin{enumerate}
 \item 任意の coherent 対象と任意の下限 $i_0$ に対し、
 その対象は $i\geq i_0$ を満たす $E_i$ の有限直和の商である。
 また任意の coherent 対象間の全射 $M\to N$ に対し、
 $\operatorname{Hom}(E_i,M)\to\operatorname{Hom}(E_i,N)$ は $i\gg0$ で全射である。
 \item Veronese 関手は基底変更に自然な同値
 $\operatorname{QGr}(\mathcal A_S)\simeq\operatorname{QGr}(\mathcal A_S^{(2)})$ を与える。
 偶数列の section 代数は前定理によって $\mathcal A_S^{(2)}$ となる。
 \item 同じ始域 $\operatorname{QGr}(\mathcal A_S)$ と終域を持つ
 二つの完全同値の偶数列上の制限の間に、全ての Hom と合成に両立する同型が与えられたなら、
 その同型は圏全体の自然同型へ一意的に延びる。
\end{enumerate}
(2) は任意の可換 $P$-代数 $S$ について成り立ち、noetherian 性を必要としない。
\end{proposition}

\begin{proof}
$\ell$ を偶数に選んだので、次数捩れ $(\ell)$ は
有限代数層模型において $\mathcal O_{Y_S}(1)$ との tensor に一致する。
通常のSerreの生成性と消滅\cite[Lemma~30.14.2]{StacksProj}を、coherent な $\mathscr A_S$-加群の台となる
$\mathcal O_{Y_S}$-加群に適用する。
例えばその大きな twist の $\mathcal O_{Y_S}$-加群としての有限生成切断は、
$\mathscr A_S$-加群としての $\mathscr A_S$ の有限直和からの全射を与える。
有限個の剰余類 $2i\bmod\ell$ について同じ議論を行えば (1) を得る。

$\ell$ は二で割れるので、$\mathcal A_S$ と $\mathcal A_S^{(2)}$ は同じ中心標準開上で
同じ次数零局所化代数を持つ。
前命題の局所化と貼合せによる構成は、この二つを同じ $\mathscr A_S$-加群圏へ送る。
これが (2) の同値であり、基底変更にも自然である。
この同値は$\pi\mathcal A_S(-2i)$を$\pi\mathcal A_S^{(2)}(-i)$へ送る。
従って商関手$\operatorname{Gr}(\mathcal A_S^{(2)})\to
\operatorname{QGr}(\mathcal A_S^{(2)})$の右随伴の非負次数部分は，
この同値の下で$\bigoplus_{i\ge0}\operatorname{Hom}(E_i,-)$に対応する。

(1) により coherent 対象 $M$ は、$E_i$ の有限直和を項とする表示
$L_1\to L_0\to M\to0$ を持つ。
二つの完全同値を $F,G$ と書く。
偶数列上の与えられた同型を有限直和に延ばし、
表示の余核を取ると $\alpha_M:F(M)\to G(M)$ が得られる。
固定した全射 $L_0\to M$ に対して、余核上の写像は
$\alpha_M F(L_0\to M)=G(L_0\to M)\alpha_{L_0}$ という等式で一意に決まり、
核を覆う $L_1$ の選択には依存しない。
二つの全射 $L_0\to M$、$L'_0\to M$ を比較するときは、
共通の全射 $L_0\oplus L'_0\to M$ と、その coherent な核の有限表示を取る。
この共通表示から得る写像は、各直和因子について上記の等式を満たす。
$F(L_0)\to F(M)$ と $F(L'_0)\to F(M)$ の全射性により、
元の二つの表示から得る写像はともにこの写像に等しい。

射 $h:M\to N$ に対しては、全射 $L_0\to M$ と $P_0\to N$ を取り、
$(h\circ(L_0\to M),P_0\to N):L_0\oplus P_0\to N$ という全射を用いる。
これから構成した $\alpha_N$ の等式を $L_0$ に制限し、
$F(L_0)\to F(M)$ が全射であることを使うと
$\alpha_NF(h)=G(h)\alpha_M$ が得られる。
この議論は $E_i$ が射影的であることや、表示への任意の射の持上げを仮定しない。
従って coherent 部分圏で表示によらない自然同型を得る。
余核表示はその一意性も与える。

最後に、locally noetherian Grothendieck 圏の各対象は
noetherian 部分対象の filtered colimit であり、
完全同値は noetherian 対象と全ての colimit を保つ。
従って coherent 部分圏上の自然同型は圏全体へ一意的に延びる。
\end{proof}

\subsection{局所 perfect 対象について得られる相対 RHom}

$\operatorname{Perf}(Y_S,\mathscr A_S)$ を、中心標準開上で
有限生成射影 $\mathscr A_S$-加群の有界複体に同型な対象の dg 圏とする。
命題\ref{prop:central-model-qgr}と
\cite[Definitions~3.1, 3.13, Theorem~3.14]{BDGMinors}により、
任意の可換 $P$-代数 $S$ について
\[
 D(\operatorname{QGr}(\mathcal A_S))^c
 \simeq \operatorname{Perf}(Y_S,\mathscr A_S)
\]
である。ここで上添字 $c$ は非有界導来圏の compact 対象を表す。
この同定に必要なのは $Y_S$ の quasi-compact 性と separated 性であり、
$S$ の noetherian 性や fiber の正則性ではない。
全係数基底上では，これを「有界coherentかつ底$S$上有限Tor振幅」という
内在的な部分圏と同一視しない。後の命題\ref{prop:central-model-intrinsic-perfect-open}は，
一般ファイバーの正則性の下で，非空基底開上に限ってその同一性を証明する。

\begin{proposition}[局所 perfect 圏の proper 性と基底変更]
\label{prop:central-model-perfect-rhom}
任意の noetherian $P$-代数 $S$ について、全ての $\mathcal O(m)$ は
$\operatorname{Perf}(Y_S,\mathscr A_S)$ に入る。
この圏の二対象 $K,L$ の相対 RHom は $S$-perfect である。
任意の noetherian $S$-代数 $S'$ に対して、標準写像
\[
 \operatorname{RHom}_{\mathscr A_S}(K,L)\otimes_S^{\mathbf L}S'
 \longrightarrow
 \operatorname{RHom}_{\mathscr A_{S'}}(K_{S'},L_{S'})
\]
は同型である。ここで右辺への移行は導来基底変更である。
正確には、中心 affine chart の代数 $\Lambda$ を使って
$K_{S'}=K\otimes_{\Lambda}^{\mathbf L}(\Lambda\otimes_SS')$
と定め、これらを貼り合わせる。
これは代数の係数拡大であり、$Y_S\to\operatorname{Spec}S$ の平坦性を仮定しないまま
台の $\mathcal O_{Y_S}$-複体の通常の導来逆像と同一視しない。
\end{proposition}

\begin{proof}
中心標準開上で $\mathscr A_S$ は
$\Lambda=(\mathcal A_S[c^{-1}])_0$ で表される。
この代数は $S$ 上平坦である。
実際、次数零局所化は有限自由 $S$-加群
$(\mathcal A_S)_{j\deg c}$ の filtered colimit である。
強次数付き性から $(\mathcal A_S[c^{-1}])_m$ は有限生成射影 $\Lambda$-加群であり、
従って全ての $\mathcal O(m)$ は局所 perfect である。

$K,L$ を局所 perfect とする。
$\mathbf R\mathcal Hom_{\mathscr A_S}(K,L)$ は有界 coherent な
$\mathcal O_{Y_S}$-複体である。
上記の各標準開で有限射影複体を用いて計算する。
有限射影右$\Lambda$-加群間のHomは，$S$-加群として$\Lambda$の有限直和の直和因子であり，
従ってこの局所Hom複体は$S$-平坦加群の有界複体となる。
このHomを自然な右$\Lambda$-加群とみなす必要はない。
有限個の開を選べるので、底に関する Tor 振幅は一様に有限である。
$Y_S\to\operatorname{Spec}S$ は射影的なので、
proper pushforward の perfect 性
\cite[Lemma 36.27.1, Tag~08EV]{StacksPerfect}
をこの $\mathcal O_{Y_S}$-複体に適用できる。
従ってその大域 RHom は $S$-perfect である。

基底変更は、中心標準開とその交叉による有限\v{C}ech 複体で確認する。
各交叉は affine であり、局所 RHom は有限射影 $\Lambda$-複体で計算できる。
$\Lambda$ の $S$-平坦性と有限射影性により、この局所計算は導来基底変更と可換する。
有限\v{C}ech totalization も導来 tensor と可換するので、表示した同型を得る。
この議論は $Y_S$ 自体が $S$ 上平坦であるとは仮定していない。
\end{proof}

\subsection{非空開上の相対対角と内在的なperfect部分圏}

全係数基底上では局所perfect圏だけを用いた。
ここでは一般の幾何学的ファイバーのtails大域次元を入力として，
有限な中心標準開被覆の各相対対角をperfectにする非空基底開を得る。
$u\ne0$と$v\ne0$の二chartだけでは$u=v=0$の境界を覆わないため，
射影中心模型の全標準開被覆を用いる。

\begin{lemma}[中心 affine 開への Ext の局所化]
\label{lem:central-affine-ext-localization}
$k$ を体、$Y$ を noetherian 射影 $k$-スキーム、
$\mathscr A$ を有限な $\mathcal O_Y$-代数とする。
豊富な可逆層 $\mathcal O_Y(1)$ の切断 $c$ に対し
$U=Y\setminus V(c)$ が affine であるとし、$j:U\hookrightarrow Y$、
$\Lambda=\Gamma(U,\mathscr A)$ と置く。
coherent な右 $\mathscr A$-加群 $M,N$ について、自然な同型
\[
 \operatorname{Ext}^i_\Lambda(M|_U,N|_U)
 \simeq\varinjlim_n
 \operatorname{Ext}^i_{\mathscr A}(M,N\otimes\mathcal O_Y(n))
 \qquad(i\geq0)
\]
がある。右辺の遷移写像は $c$ による乗法である。
さらに $\operatorname{coh}(Y,\mathscr A)$ の大域次元が高々 $d$ なら、
$\operatorname{gldim}\Lambda\leq d$ である。
\end{lemma}

\begin{proof}
まず coherent 対象 $M$ に対して、Serre の生成性を繰り返すと、
各項が $\mathscr A\otimes\mathcal O_Y(-n)$ の有限直和である
有界上方の分解 $L^\bullet\to M$ を取れる。
核は常に coherent なので、この操作を任意の長さまで継続できる。
$L^\bullet$ の各項は局所的に有限自由な $\mathscr A$-加群である。
$Y$ の有限 affine 開被覆を選ぶと、
$\mathcal Hom_{\mathscr A}(L^\bullet,N)$ の有限\v{C}ech totalization は
$\operatorname{RHom}_{\mathscr A}(M,N)$ を計算する。
sheaf Hom の次数および\v{C}ech の次数はともに非負で、後者は一様に有界である。
従って固定した total degree に関わる項は有限個である。
局所有限自由性、affine 上の切断、有限\v{C}ech 複体は filtered colimit と可換するので、
$\operatorname{Ext}^i_{\mathscr A}(M,-)$ は quasi-coherent 加群の filtered colimit と可換する。

可逆層を局所的に自明化して通常の環の局所化を使えば
\[
 j_*j^*N=\varinjlim_n N\otimes\mathcal O_Y(n)
\]
となる。
$Y$ は separated で $U$ は affine なので $j$ は affine である。
従って quasi-coherent 加群に対する $j_*$ は完全であり、
$j^*$ との導来随伴と直前の colimit との可換性から表示した Ext 同型を得る。
ここでcoherent対象間のExtはcoherent圏でもquasi-coherent圏でも同じである。
実際，locally noetherian Grothendieck圏ではnoetherian対象への任意の全射は，
あるnoetherian部分対象上で既に全射となる。
従って\cite[Lemma~13.17.4, Tag~0FCL]{StacksAbel}をnoetherian対象のSerre部分圏に
適用でき，必要な有界上方導来圏の比較を得る。

任意の有限生成 $\Lambda$-加群 $Q$ は coherent な $\mathscr A$-加群へ延長できる。
実際、$Q$ の有限個の生成元は
$\Gamma(U,Q)=\Gamma(Y,j_*Q)$ の切断であり、
対応する $\mathscr A^{\oplus r}\to j_*Q$ の像は coherent で、$U$ 上では $Q$ に一致する。
従って大域次元の仮定と Ext の同型により、
全ての有限生成 $\Lambda$-加群 $Q,Q'$ に対して
$\operatorname{Ext}^{i}_\Lambda(Q,Q')=0$、$i>d$ となる。
$\Lambda$ は affine noetherian 可換環上有限なので左右 noetherian である。
有限生成 $Q$ の有限自由分解における $d$ 次 syzygy は有限生成であり、
それに対する有限自由全射の核も有限生成である。
その短完全列は上記の Ext 消滅によって分裂するので、$\operatorname{pd}_\Lambda Q\leq d$ である。
とくに全ての巡回加群 $\Lambda/I$ の射影次元が高々 $d$ である。
Baer の判定を dimension shifting とともに使えば、全ての加群の入射次元が高々 $d$ となり、
$\operatorname{gldim}\Lambda\leq d$ を得る。
\end{proof}

\begin{lemma}[相対対角と有限 Tor 振幅]
\label{lem:relative-diagonal-finite-tor}
$S$ を可換 noetherian 環、$\Lambda$ を $S$ 上平坦な左右 noetherian 代数とする。
$\Lambda$ が対角両側加群として
$\Lambda^{\mathrm{op}}\otimes_S\Lambda$ 上 perfect であると仮定する。
このとき、$K\in D^b_{\mathrm{fg}}(\Lambda)$ に対し
\[
 K\in\operatorname{Perf}(\Lambda)
 \quad\Longleftrightarrow\quad
 K\text{ は }S\text{ 上有限 Tor 振幅を持つ}.
\]
\end{lemma}

\begin{proof}
左辺から右辺は $\Lambda$ の $S$-平坦性から従う。
逆に $K$ の $S$-Tor 振幅を $[a,b]$ とする。
対角 $\Lambda$ の有界有限射影
$\Lambda^{\mathrm{op}}\otimes_S\Lambda$-分解を $Q^\bullet$ とする。
任意の左 $\Lambda$-加群 $N$ に対し
\[
 K\otimes_\Lambda^{\mathbf L}N
 \simeq K\otimes_\Lambda^{\mathbf L}Q^\bullet
              \otimes_\Lambda^{\mathbf L}N
\]
を使う。
$Q^j$ は有限個の外積 $\Lambda\otimes_S\Lambda$ の直和因子なので、
各項の derived tensor は有限個の
$K\otimes_S^{\mathbf L}N$ の直和因子である。
$Q^\bullet$ の長さだけ Tor 振幅を広げれば、全ての $N$ に共通する有界区間を得る。
従って $K$ は $\Lambda$ 上有限 flat dimension を持つ。
$\Lambda$ は noetherian で $K$ の cohomology は有界有限生成なので、
有限生成自由加群の有界上方分解を十分低い次数で切ると、最後の syzygy は有限表示かつ flat になる。
これは射影的であり、$K$ は perfect である。
\end{proof}

この補題は有限次元 $S$-加群性を仮定しない。
$K\otimes_S\Lambda$ は一般には有限生成射影 $\Lambda$-加群ではないので、
それを誤って有限射影と扱うことは避ける。
証明でまず得るのは有限 flat dimension であり、
coherent 性と noetherian 性によって最後に有限射影分解へ移る。

\begin{lemma}[中心上有限な affine 族の一般点からの拡張]
\label{lem:relative-diagonal-generic-spread}
$P$ を noetherian 整域、$K=\operatorname{Frac}P$ を標数零の体とする。
$D$ を有限型可換 $P$-代数、$\Lambda$ を有限な $D$-代数とし、
$\Lambda$ は $P$ 上平坦であるとする。
一般 fiber の対角 $\Lambda_K$ が
$\Lambda_K^{\mathrm{op}}\otimes_K\Lambda_K$ 上 perfect であると仮定する。
この仮定は、例えば $\Lambda_K$ または $\Lambda_{\overline K}$ が有限大域次元を持てば満たされる。
このとき
ある $0\ne p\in P$ に対し $\Lambda_p$ の相対対角は
$\Lambda_p^{\mathrm{op}}\otimes_{P_p}\Lambda_p$ 上 perfect である。
その有限射影分解は任意の $P_p$-基底変更へ移せる。
\end{lemma}

\begin{proof}
有限大域次元から表示した generic perfect 性が従うという補足を先に確認する。
Rennemo--Segal--Van den Bergh\cite[Proposition~2.1, Lemmas~3.3--3.6, Corollary~3.9]{RSVdB}の証明の
steps (3)--(5) は、標数零の体上で有限型可換代数上有限な代数について、
有限大域次元から homological smoothness を導く。
ここでは彼らの Lemma 3.4（反対代数）、Lemma 3.5（PI 代数の tensor）、
Lemma 3.6（代数閉包からの降下）が使われる。
幾何学的一般 fiber $\Lambda_{\overline K}$ の有限大域次元だけを入力とする場合は、
まず代数閉体上の steps (4) によりその対角を perfect とし、
Lemma 3.6 により $K$ 上へ降下すればよい。

包絡代数 $\Lambda^{\mathrm{op}}\otimes_P\Lambda$ は
有限型可換 $P$-代数 $D\otimes_PD$ 上有限であり、
対角 $\Lambda$ はその有限生成加群である。
これは同論文 Corollary 3.9 の適用範囲である。
必要な spreading の議論を明示するため、包絡代数を $\Gamma$ と書き、
$\Lambda$ の有限生成自由 $\Gamma$-加群による分解を取る。
$\Gamma$ は noetherian で、$\Lambda_K$ は perfect なので、
十分高い syzygy $M$ は有限表示であり、$M_K$ は有限生成射影 $\Gamma_K$-加群である。
有限自由加群を経由する分裂
\[
 M_K\xrightarrow{i_K}\Gamma_K^{\oplus r}
       \xrightarrow{p_K}M_K,\qquad p_Ki_K=1
\]
を取る。
$M$ の有限表示性により Hom は中心局所化と可換するので、
分母を払えば、ある $0\ne s\in P$ の後で $i_K,p_K$ は $i_s,p_s$ に降りる。
$p_si_s-1$ は一般点で零であり、$M_s$ の有限個の生成元でその値の分母をさらに払うと、
ある $0\ne t\in P$ の後で $p_{st}i_{st}=1$ が成り立つ。
従って $M_{st}$ は有限生成射影であり、相対対角は $\Gamma_{st}$ 上 perfect となる。
この分裂の分母払いは既知の generic projectivity の議論であり、新規性は主張しない。
$\Lambda$ は $P$-平坦なので、対角の有界有限射影分解に任意の基底変更を適用しても、
導来基底変更した対角は通常の $\Lambda\otimes_PS$ と一致する。
包絡代数も基底変更した包絡代数に一致するため、所望の分解を得る。
\end{proof}

\begin{proposition}[非空開上の内在的な有限 Tor 部分圏]
\label{prop:central-model-intrinsic-perfect-open}
定理\ref{thm:central-model-finite}の仮定を満たす四係数族を考え、
$K=\operatorname{Frac}P$ とする。
追加の入力として
\begin{equation}\label{eq:geometric-generic-tails-input}
 \operatorname{gldim}\operatorname{qgr}\mathcal A_{\overline K}<\infty
\end{equation}
を仮定する。
ある非零 $p\in P$ が存在し、任意の noetherian $P_p$-代数 $S$ について
\[
 \operatorname{Perf}(Y_S,\mathscr A_S)
 =\{K'\in D^b_{\mathrm{coh}}(Y_S,\mathscr A_S):
          K'\text{ は }S\text{ 上有限 Tor 振幅を持つ}\}
\]
となる。
この部分圏は代数の係数拡大による導来基底変更に保たれ、
二対象間の相対 RHom は $S$-perfect であり、任意の noetherian 基底変更と可換する。
\end{proposition}

\begin{proof}
$D=C^{(\ell)}$ の有限個の次数一元 $c_i$ により、$Y$ の中心標準開被覆を取る。
その代数を
\[
 D_i=(D[c_i^{-1}])_0,\qquad
 \Lambda_i=(\mathcal A[c_i^{-1}])_0
\]
と書く。
$D_i$ は有限型可換 $P$-代数であり、$\Lambda_i$ は $D_i$ 上有限で $P$-平坦である。
後者の平坦性は次数零局所化が有限自由 $P$-加群の filtered colimit であることによる。

仮定 \eqref{eq:geometric-generic-tails-input} と中心模型の同値、
補題 \ref{lem:central-affine-ext-localization} を
$Y_{\overline K}$ の各中心 affine 開に適用する。
各 $\Lambda_{i,\overline K}$ は有限大域次元を持つ。
従って前補題の議論によって $\Lambda_{i,K}$ の対角は perfect である。
さらに同補題により、ある非零 $p_i\in P$ の後で $\Lambda_{i,p_i}$ の相対対角は perfect となる。
$p=\prod_i p_i$ とすれば有限個の chart で同時にこれが成立する。
この段階で重要なのは $u,v$ の二 chart ではなく、$Y$ 全体を覆う有限中心被覆を用いたことである。

任意の noetherian $P_p$-代数 $S$ に対し、
$\Lambda_{i,S}$ は $S$-平坦で左右 noetherian であり、相対対角の有限射影分解を持つ。
補題 \ref{lem:relative-diagonal-finite-tor} によって、各 chart で
有界 coherent かつ有限 $S$-Tor 振幅を持つ対象は $\Lambda_{i,S}$-perfect に一致する。
内在的な Tor 振幅の定義は $S$-加群作用から与えられるので、この chart への制限と可換する。
被覆が有限であるため、局所的な有限 Tor 振幅は共通の有限区間で抑えられる。
従って表示した大域的な部分圏の同一性を得る。
基底変更と相対 RHom に関する主張は、有限中心模型における
命題\ref{prop:central-model-perfect-rhom}から従う。
\end{proof}

\paragraph{具体的な族への適用入力と残る問題.}
$B_*$ の固定四係数族については、既存の開な正則性の定理
\ref{CSBGthm:global-open} が、PBW/Koszul 解消と AS 条件の体拡大への両立の下で
任意の代数閉拡大体の fiber の正則性を与える。
これを非零多項式の開集合の geometric generic point へ適用すれば
式 \eqref{eq:geometric-generic-tails-input} の値は三となる。
ここで $K$ 自身への tails 大域次元の降下を省略して仮定したわけではない。
局所化の補題は最初に $\overline K$ 上で適用し、
affine 対角の perfect 性を RSV の Lemma 3.6 により $K$ へ降下している。

条件\textup{(C1),(C2)}は上の全四係数族で検証された。
この命題を条件\textup{(C1)--(C4)}の比較へ用いる場合、さらに得られるのは
非空基底開上の (C3) の内在的部分圏の同定と、その部分圏上の
perfect な相対 RHom および基底変更である。
これは(C4)のRHom部分を満たす。
相対Serre関手の存在と$(-2)[3]$という同定は，
第\ref{sec:relative-serre-projective}節で全四係数基底上に証明する。
そこでcompact対象を使えば，Hom比較には本節の非空開への制限が不要となる。
大域 dg smooth 性には、局所対角に加えて大域的な生成と Morita 比較が必要である。
第\ref{sec:dg-smooth-torsion-families}節では、左右の正則性の下で
外積 classical generator と対角 kernel による比較を与え、
有限位数境界上の相対 smooth proper 開を別に構成する。

